%
% architecture.tex
%
% Vox Technical Architecture Specification
% Written for principal engineers evaluating the system.
%
% Compile: pdflatex architecture.tex
%

\documentclass[a4paper,10pt]{article}
\usepackage[margin=1in]{geometry}
\usepackage{amssymb}
\usepackage{booktabs}
\usepackage{listings}
\usepackage{xcolor}
\usepackage{enumitem}
\usepackage{longtable}
\usepackage[colorlinks=true,linkcolor=blue,citecolor=blue,urlcolor=blue]{hyperref}

\lstset{
  basicstyle=\ttfamily\small,
  keywordstyle=\color{blue!80!black},
  commentstyle=\color{green!50!black},
  stringstyle=\color{red!60!black},
  breaklines=true,
  frame=single,
  framerule=0.3pt,
  rulecolor=\color{gray!50},
  backgroundcolor=\color{gray!5},
  columns=fullflexible,
}

\newcommand{\modpath}[1]{\texttt{src/punt\_vox/#1}}
\newcommand{\tool}[1]{\texttt{#1}}
\newcommand{\field}[1]{\texttt{#1}}
\newcommand{\hook}[1]{\texttt{#1}}
\newcommand{\cmd}[1]{\texttt{/#1}}

\begin{document}

\title{Vox: Technical Architecture Specification}
\author{Text-to-Speech Engine for AI Agents}
\date{July 2026 --- v4.11.0}
\hypersetup{
  pdftitle={Vox: Technical Architecture Specification},
  pdfauthor={Punt Labs},
  pdfsubject={Architecture},
  pdfcreator={pdflatex}
}
\maketitle

\tableofcontents
\newpage

% ===================================================================
\section{Overview}
% ===================================================================

Vox is a text-to-speech engine with three projection surfaces from a
single codebase: Python CLI, MCP server, and Claude Code plugin.  It
supports five synthesis backends (ElevenLabs, OpenAI, AWS Polly, macOS
\texttt{say}, Linux \texttt{espeak-ng}) behind a provider-agnostic
orchestration layer.  Audio playback is serialized via file locking.
Session state lives in a YAML frontmatter config file.

The system does \emph{not} stream audio.  Every synthesis call produces
a complete MP3 file before playback begins.  There is no audio mixing and
no real-time pitch manipulation.  Vox is a batch synthesizer with
a playback queue.  The \texttt{voxd} daemon accepts requests over
WebSocket and serializes playback.  Since v4.x it also generates,
persists, and plays background music as an ownership-free
\emph{Program} --- a named on-disk pool of instrumental tracks
(\texttt{elevenlabs\_music}) that fills in the background, auto-advances,
and shuffle-rotates once full --- through a separate reduced-volume
subprocess outside the speech queue, and can be driven from a remote
host over the network (\texttt{VOXD\_HOST} / \texttt{VOXD\_PORT} /
\texttt{VOXD\_TOKEN}).

\subsection{Design Principles}

\begin{enumerate}[label=\textbf{P\arabic*}.]
  \item \textbf{Provider-agnostic orchestration.}  The core module
    knows nothing about any specific TTS API.  Providers implement
    a \texttt{TTSProvider} protocol; the core synthesizes, splits,
    stitches, and plays without knowing which backend produced the
    audio.
  \item \textbf{No audio until opt-in.}  Vox is silent by default.
    The \field{notify} config field defaults to \texttt{"n"}.  No
    hooks produce sound, no chimes play, and no synthesis calls fire
    until the user explicitly enables notifications via \cmd{unmute}
    or \cmd{vox}.
  \item \textbf{Session-scoped state.}  Mutable state (voice,
    provider, mood) lives in the MCP server's memory, one instance
    per Claude Code session.  Initial state is seeded from the
    per-repo split config (\texttt{.punt-labs/vox/vox.md} for durable
    preferences, \texttt{.punt-labs/vox/vox.local.md} for ephemeral
    session state) at startup; hook handlers also read and write these
    files for signal accumulation.  The \texttt{voxd} daemon
    is stateless with respect to \emph{speech} sessions---it synthesizes
    and plays whatever clients send it.  Background music is the one
    piece of daemon-owned state: a single persisted, ownership-free
    \emph{Program} that any session drives and no session owns
    (\S\ref{sec:music}, DES-041).
  \item \textbf{Fail fast.}  Invalid input raises exceptions.
    No defensive fallbacks, no silent degradation.  If a voice
    doesn't exist, the call fails with a
    \texttt{VoiceNotFoundError}---it does not silently substitute
    another voice.
  \item \textbf{Serialized playback.}  Concurrent synthesis calls
    never overlap audio output.  The \texttt{voxd} daemon owns an
    in-memory playback queue with a single consumer---audio items
    play sequentially across all connected clients.
\end{enumerate}

\subsection{Component Map}

\texttt{src/punt\_vox/} holds roughly thirty modules plus three
sub-packages (\modpath{voxd/}, \modpath{service/}, \modpath{providers/}).
The daemon and the system-service installer were split from single
files into packages as they grew.  Domain types were split out of a
single \texttt{types.py} into purpose-scoped modules.  The tables below
list the top-level modules first, then the package contents.

\subsubsection*{Top-level modules}

\begin{center}
\begin{longtable}{llp{7cm}}
\toprule
\textbf{Module} & \textbf{File} & \textbf{Responsibility} \\
\midrule
\endfirsthead
\toprule
\textbf{Module} & \textbf{File} & \textbf{Responsibility} \\
\midrule
\endhead
Types      & \modpath{types.py}      & Provider protocol and shared
             domain types (\texttt{TTSProvider}, \texttt{AudioProvider},
             \texttt{MergeStrategy}) \\
Audio types & \modpath{types\_audio.py} & Audio request/result value
             objects: \texttt{AudioRequest}, \texttt{AudioResult},
             \texttt{HealthCheck} \\
Error types & \modpath{types\_errors.py} & Domain error hierarchy
             (\texttt{VoiceNotFoundError} and siblings) \\
Synthesis types & \modpath{types\_synthesis.py} & \texttt{SynthesisSpec}
             --- validated voice-settings bundle passed to the daemon \\
Core       & \modpath{core.py}       & \texttt{TTSClient}---provider-agnostic
             orchestration: text splitting, batch synthesis, pair
             stitching, audio merge \\
Config     & \modpath{config.py}     & \texttt{ConfigStore}/\texttt{VoxConfig}:
             read/write the split config (\texttt{vox.md} durable +
             \texttt{vox.local.md} ephemeral); \texttt{DURABLE\_KEYS} /
             \texttt{EPHEMERAL\_KEYS} routing \\
Dirs       & \modpath{dirs.py}       & Per-repo config-dir discovery
             (\texttt{find\_config\_dir}) and user-content dirs
             (\texttt{default\_output\_dir}, \texttt{music\_output\_dir}) \\
Paths      & \modpath{paths.py}      & \emph{Authoritative} per-user
             state paths: \texttt{\~{}/.punt-labs/vox} and its
             \texttt{logs/}, \texttt{run/}, \texttt{cache/} subdirs;
             \texttt{keys\_env\_file()}, \texttt{ensure\_user\_dirs()} \\
Resolve    & \modpath{resolve.py}    & Voice/language resolution, output
             directory resolution, vibe tag application \\
Output     & \modpath{output.py}     & Output path resolution:
             \texttt{VOX\_OUTPUT\_DIR} env, \texttt{\~{}/Music/vox}
             fallback \\
Output fmt & \modpath{output\_formatter.py} & Formats MCP tool output
             into the compact UI panel text the suppress-output hook emits \\
Playback   & \modpath{playback.py}   & \texttt{flock}-serialized playback:
             \texttt{play\_audio()} (blocking),
             \texttt{enqueue()} (detached subprocess) \\
Hooks      & \modpath{hooks.py}      & Claude Code hook dispatchers: Stop,
             PostToolUse, Notification, PreCompact,
             UserPromptSubmit, SubagentStart/Stop, SessionEnd \\
Signal     & \modpath{signal.py}     & \texttt{Signal} / \texttt{SignalLog}
             domain extracted from \texttt{hooks.py}: classification and
             timestamped accumulation \\
Hook envelope & \modpath{hook\_envelope.py} & Typed wrapper for a hook's
             JSON-in / JSON-out envelope \\
Hook payload & \modpath{hook\_payload.py} & Parses the raw hook stdin
             payload into typed fields \\
Chime      & \modpath{chime.py}      & Mood-aware chime file selection \\
Music (CLI) & \modpath{cli\_music.py}  & Client-side \cmd{music} command
             surface: on / off / play / list / next / loop and
             \texttt{playlist:N} part addressing \\
Program client & \modpath{program\_gateway.py} & Translate music/Program
             commands into \texttt{voxd} Program RPCs over the WebSocket;
             thin, no business logic (with \modpath{program\_control.py}) \\
Music prompts & \modpath{music\_prompts.py} & Scaffolds the LLM-authored
             \texttt{base\_prompt} plus twelve genre \texttt{variations}
             passed to \texttt{voxd} for track generation \\
Cache      & \modpath{cache.py}      & MP3 quip cache, content-addressed
             by (text, voice, provider) \\
Vibe       & \modpath{vibe.py}       & Vibe domain: mood text $\to$
             ElevenLabs expressive-tag resolution, extracted from
             \texttt{hooks.py} \\
Voices     & \modpath{voices.py}     & Voice personality blurbs for roster
             display \\
Quips      & \modpath{quips.py}      & Immutable phrase pools for hook
             speech (stop, permission, idle, farewell, etc.) \\
Mood       & \modpath{mood.py}       & Mood family classification
             (bright/dark/neutral) \\
Normalize  & \modpath{normalize.py}  & Text normalization for speech
             (snake\_case, camelCase, abbreviation expansion) \\
Keys       & \modpath{keys.py}       & Provider API-key presence detection
             and reporting \\
API key resolver & \modpath{api\_key\_resolver.py} & \texttt{ApiKeyResolver}
             --- resolves a per-call key from exactly one of four
             mutually-exclusive sources (\texttt{--api-key-file},
             \texttt{--api-key-stdin}, \texttt{VOX\_API\_KEY},
             \texttt{--api-key}) \\
Doctor     & \modpath{doctor.py}     & \texttt{vox doctor} health checks:
             Python, player binary, provider keys, daemon staleness \\
Daemon restarter & \modpath{daemon\_restarter.py} & Detects a stale
             running \texttt{voxd} (version mismatch) and restarts it via
             the platform service backend \\
Desktop install & \modpath{desktop\_install.py} & \texttt{DesktopInstaller}
             --- registers the \texttt{vox} MCP server in Claude Desktop's
             config \emph{without} embedding provider secrets; the daemon
             reads keys from \texttt{keys.env} at startup \\
Logging    & \modpath{logging\_config.py} & Rotating file log under
             \texttt{\~{}/.punt-labs/vox/logs/} \\
Watcher    & \modpath{watcher.py}    & \texttt{SessionWatcher}---daemon thread
             tailing Claude session JSONL files for real-time
             signal classification and autonomous chime/speech
             in continuous mode \\
Client     & \modpath{client.py}     & WebSocket client for \texttt{voxd}:
             \texttt{VoxClient} (async) / \texttt{VoxClientSync} (sync),
             split across \modpath{client\_gateway.py},
             \modpath{client\_env.py}, and \modpath{client\_errors.py}.
             Lightweight---stdlib + websockets only \\
Server     & \modpath{server.py}     & FastMCP server (key: \texttt{mic}):
             the eleven MCP tools (see \S\ref{sec:mcp}) \\
CLI        & \modpath{\_\_main\_\_.py} & Typer CLI: unmute, record, vibe,
             music, notify, speak, voice, status, doctor, version, hook,
             daemon subcommands \\
Applet     & \modpath{applet.py}     & Lux display applet: element tree
             for the visual status surface \\
\bottomrule
\end{longtable}
\end{center}

\subsubsection*{Sub-packages}

\begin{center}
\begin{longtable}{llp{7cm}}
\toprule
\textbf{Package} & \textbf{Module} & \textbf{Responsibility} \\
\midrule
\endfirsthead
\toprule
\textbf{Package} & \textbf{Module} & \textbf{Responsibility} \\
\midrule
\endhead
\modpath{voxd/} & \texttt{daemon.py} & Daemon entry point; owns
             \texttt{DEFAULT\_PORT=8421}, \texttt{DEFAULT\_HOST}, lifecycle \\
        & \texttt{router.py} & WebSocket message dispatch to handlers \\
        & \texttt{speech\_handlers.py} & synthesize / chime / record RPCs \\
        & \texttt{system\_handlers.py} & voices / health / control RPCs \\
        & \texttt{synthesis.py} & Provider selection and synthesis pipeline \\
        & \texttt{playback.py} & Daemon-side playback queue and consumer \\
        & \texttt{dedup.py} & Content-addressed replay dedup with TTL \\
        & \texttt{health.py} & Health-response assembly \\
        & \texttt{chimes.py} & Chime asset resolution inside the daemon \\
        & \texttt{config.py} & Daemon runtime config; \texttt{load\_keys()}
             reads \texttt{keys.env} into the environment at startup \\
        & \texttt{synthesis\_result.py} & Typed synthesis outcome carried
             back across the WebSocket boundary \\
        & \texttt{programs/} & Background-music \emph{Program} package
             (ownership-free, persisted; \S\ref{sec:music}):
             \texttt{program.py} / \texttt{state.py} / \texttt{invariants.py}
             (the state machine and its sixteen invariants);
             \texttt{control\_channel.py} plus the typed
             \texttt{*\_signal.py} (single-writer control);
             \texttt{filler.py} / \texttt{fill\_reconciler.py} /
             \texttt{producer.py} / \texttt{retry\_machine.py} (background
             fill and resilient retry); \texttt{rotate\_policy.py} /
             \texttt{player.py} / \texttt{subprocess\_player.py}
             (auto-advance, shuffle-rotation, reduced-volume playback);
             \texttt{filesystem\_store.py} / \texttt{manifest.py} /
             \texttt{part\_tags.py} (persistence and ID3 tagging); the
             per-command handlers; and \texttt{service.py} (composition
             root) \\
\midrule
\modpath{service/} & \texttt{installer.py} & \texttt{ServiceInstaller}
             composing the platform backends; \texttt{install} /
             \texttt{uninstall} / health verification \\
        & \texttt{launchd.py} & macOS \texttt{LaunchdBackend} ---
             user LaunchAgent under \texttt{\~{}/Library/LaunchAgents} \\
        & \texttt{launchctl.py} & \texttt{LaunchctlAgent} --- race-free
             launchd domain control (bootout-and-wait, then bootstrap with
             retry-on-error-5), shared by install and restart \\
        & \texttt{systemd.py} & Linux \texttt{SystemdBackend} ---
             unit under \texttt{/etc/systemd/system} \\
        & \texttt{process.py} & \texttt{ProcessManager}: port-free
             enforcement, stale-daemon kill \\
        & \texttt{keys\_env.py} & \texttt{KeysEnvWriter}: writes provider
             keys to \texttt{keys.env} \\
\midrule
\modpath{providers/} & \texttt{elevenlabs.py} & ElevenLabs TTS backend \\
        & \texttt{elevenlabs\_music.py} & ElevenLabs music-generation backend \\
        & \texttt{openai.py} & OpenAI TTS backend \\
        & \texttt{polly.py} & AWS Polly backend \\
        & \texttt{say.py} & macOS \texttt{say} fallback \\
        & \texttt{espeak.py} & Linux \texttt{espeak-ng} fallback \\
        & \texttt{local\_play.py} & Local-only playback provider \\
        & \texttt{chunked.py} & Shared chunk-and-stitch helper \\
        & \texttt{convert.py} & Format conversion (AIFF/WAV $\to$ MP3) \\
        & \texttt{voice\_resolver.py} & Shared voice-name resolution \\
\bottomrule
\end{longtable}
\end{center}


% ===================================================================
\section{System Architecture}
% ===================================================================

\subsection{Daemon and Connectivity}

The system has two binaries from one package (\texttt{punt-vox}):

\begin{description}
  \item[\texttt{voxd}] System-level audio daemon.  Listens on port
    8421 via WebSocket.  Owns synthesis (providers), playback queue,
    audio deduplication, and synthesis cache.  Knows nothing about
    MCP, hooks, projects, or Claude Code.
  \item[\texttt{vox}] Everything else: CLI commands, MCP server
    (\texttt{vox mcp}), hook handlers (\texttt{vox hook <event>}).
    All are WebSocket clients of \texttt{voxd}.
\end{description}

\begin{verbatim}
Claude Code <-- stdio --> vox mcp -- WebSocket --> voxd :8421
Hook scripts --> vox hook <event> -- WebSocket -->    |
Shell        --> vox unmute "hi"  -- WebSocket -->    |
                                                      v
                                                   speakers
\end{verbatim}

The MCP session (Claude~Code~$\leftrightarrow$~\texttt{vox mcp}) is
stdio---unaffected by daemon restarts.  The WebSocket connection
(\texttt{vox mcp}~$\leftrightarrow$~\texttt{voxd}) reconnects
automatically.

\subsection{Remote voxd}

By default \texttt{voxd} binds \texttt{127.0.0.1} and clients discover
its port and token from the local \texttt{run/serve.\{port,token\}}
files.  DES-037 lets the client half run on a different machine from the
daemon --- the common case being SSH from a laptop with speakers to a
headless server where synthesis and playback would otherwise be
inaudible.  Four environment variables configure it:

\begin{center}
\begin{tabular}{lll}
\toprule
\textbf{Var} & \textbf{Side} & \textbf{Default} \\
\midrule
\texttt{VOXD\_HOST}  & client & \texttt{127.0.0.1} \\
\texttt{VOXD\_PORT}  & client & from \texttt{serve.port} \\
\texttt{VOXD\_TOKEN} & client & from \texttt{serve.token} \\
\texttt{VOXD\_BIND}  & server & \texttt{127.0.0.1} \\
\bottomrule
\end{tabular}
\end{center}

Resolution is \texttt{explicit arg > env var > file > default}.  Two
deployment models are supported: direct network (same LAN) and SSH
tunnel (different networks).  The token is the security boundary and is
redacted from access logs; there is no TLS on \texttt{voxd} itself.  See
\texttt{docs/remote-setup.md}.

\subsection{System Paths}

\texttt{voxd} runs as a single user, so \emph{all} of its state lives
under one per-user root: \texttt{\~{}/.punt-labs/vox/}.  The same scheme
applies on macOS and Linux --- there is no FHS split and no Homebrew
prefix.  \modpath{paths.py} is the single source of truth for these
locations; \texttt{ensure\_user\_dirs()} creates the root and its
\texttt{logs/}, \texttt{run/}, and \texttt{cache/} subdirs \texttt{chmod
0700} because each holds private per-user state (API keys, spoken-text
logs, the auth token, cached audio).

An earlier release resolved these to FHS system paths
(\texttt{/etc/vox}, \texttt{/var/log/vox}, \texttt{/var/run/vox}) on
Linux and Homebrew-prefix paths on macOS.  That was a regression: it
stranded user API keys on upgrade, required \texttt{sudo} to edit
personal tokens, and created a chown mismatch where root owned the file
\texttt{voxd} was told to read.  It has been removed.

\begin{center}
\begin{tabular}{lll}
\toprule
\textbf{Purpose} & \textbf{macOS} & \textbf{Linux} \\
\midrule
Keys     & \multicolumn{2}{l}{\texttt{\~{}/.punt-labs/vox/keys.env}} \\
Cache    & \multicolumn{2}{l}{\texttt{\~{}/.punt-labs/vox/cache/}} \\
Logs     & \multicolumn{2}{l}{\texttt{\~{}/.punt-labs/vox/logs/voxd.log}} \\
Runtime  & \multicolumn{2}{l}{\texttt{\~{}/.punt-labs/vox/run/serve.\{port,token\}}} \\
Service  & \texttt{\~{}/Library/LaunchAgents/}
         & \texttt{/etc/systemd/system/} \\
         & \texttt{com.punt-labs.voxd.plist}
         & \texttt{voxd.service} \\
\bottomrule
\end{tabular}
\end{center}

The \texttt{run/} directory is created \texttt{chmod 0700} so other
local users cannot read the auth token.  Per-project enablement is a
separate concern: the split config
(\texttt{.punt-labs/vox/vox.md} + \texttt{vox.local.md}) is a
client-side pair in the project directory.  The daemon never reads it.

\subsection{Service Install and Identity}

\begin{verbatim}
vox daemon install
\end{verbatim}

Runs as the invoking user.  \texttt{install()} refuses to run with
\texttt{os.geteuid() == 0}, so users who run \texttt{sudo vox daemon
install} out of habit get an explicit error directing them to
re-run without sudo.  Per-user state under
\texttt{\textasciitilde/.punt-labs/vox/} (keys, logs, runtime state,
cache) is created with normal user permissions---no chown, no
symlink defenses.  \texttt{voxd} runs as the \emph{installing user},
not root: it needs audio device access tied to the desktop session
(CoreAudio on macOS, PulseAudio/PipeWire on Linux).

The two platforms differ in whether \texttt{sudo} is needed at all:

\begin{description}
  \item[macOS --- zero sudo (DES-038).] Since v4.9.0 the daemon is a
    user \emph{LaunchAgent} at
    \texttt{\~{}/Library/LaunchAgents/com.punt-labs.voxd.plist}, not a
    root-owned LaunchDaemon.  Both \texttt{install} and \texttt{uninstall}
    are fully sudo-free: \texttt{mkdir -p \~{}/Library/LaunchAgents},
    write the plist, free the port, then \texttt{launchctl bootstrap
    gui/<uid>} and \texttt{launchctl kickstart}.  The move fixed a 7x
    synthesis slowdown --- LaunchDaemons run at background QoS with CPU
    demotion and I/O deprioritization; LaunchAgents run at user QoS.
    The plist drops \texttt{UserName} (invalid for agents) and adds
    \texttt{ProcessType=Interactive} to prevent App~Nap throttling on
    the windowless daemon.  There is no LaunchDaemon-to-LaunchAgent
    migration path --- it was removed in v4.9.0 rather than shipped
    (only a handful of installs existed, and a compat shim would violate
    the no-backwards-compat-shim rule).
  \item[Linux --- sudo for the unit file only.] The systemd unit at
    \texttt{/etc/systemd/system/voxd.service} lives in a root-owned
    system directory, so placing it, reloading the manager, enabling it
    at boot, and cycling the process still escalate through
    \texttt{sudo}.  The unit sets \texttt{User=} to the installing user
    so the daemon still runs unprivileged.
\end{description}

After registration, \texttt{install()} polls \texttt{voxd}'s health
endpoint until it answers (5\,s deadline) --- \texttt{launchctl} /
\texttt{systemctl} registration proves only that the job is scheduled,
not that the daemon bound its port and stayed up.  See DES-038 (and the
superseded DES-029) in \texttt{DESIGN.md} for the full rationale.

\subsection{Data Paths}

Three primary data paths define the system's runtime behavior.

\subsection{TTS Request Path (Text $\to$ Audio File $\to$ Playback)}

\begin{enumerate}
  \item Agent calls MCP tool (e.g., \tool{unmute(text="Hello")}).
  \item Server resolves provider: explicit parameter $\to$ session
    config $\to$ auto-detection (environment variable $\to$ API key
    probe $\to$ platform fallback).
  \item Server resolves voice: explicit parameter $\to$ session config
    $\to$ language-aware default $\to$ provider default.
  \item Server builds \texttt{AudioRequest} list (one per segment).
  \item Server predicts output paths and returns immediately to the
    agent with predicted \texttt{AudioResult} metadata.
  \item Background daemon thread: \texttt{TTSClient.synthesize\_batch()}
    calls the provider's \texttt{synthesize()} for each request.
  \item If multiple segments: \texttt{stitch\_audio()} concatenates
    MP3 files with configurable inter-segment silence (default
    500\,ms) plus 150\,ms trailing silence.
  \item \texttt{enqueue(path)} copies the file to
    \texttt{\~{}/.punt-labs/vox/pending/}, spawns a detached subprocess.
  \item Subprocess acquires \texttt{flock(LOCK\_EX)} on
    \texttt{\~{}/.punt-labs/vox/playback.lock}, plays via
    \texttt{afplay} (macOS) or \texttt{ffplay} (cross-platform),
    deletes the pending copy on completion.
\end{enumerate}

The agent receives the MCP response at step~5.  Steps~6--9 happen
asynchronously.  The agent never blocks on synthesis or playback.

\subsection{Notification Path (Hook Event $\to$ Audio)}

\begin{enumerate}
  \item Claude Code fires a lifecycle hook (Stop, Notification,
    PreCompact, etc.).
  \item Shell script dispatcher (e.g., \texttt{hooks/notify.sh})
    invokes \texttt{vox hook stop} with JSON on stdin.
  \item Python hook handler reads session config, checks
    \field{notify} and \field{speak} fields.
  \item If \field{notify} is \texttt{"n"}: no-op, immediate return.
  \item If \field{speak} is \texttt{"n"} (chime mode): resolves
    mood-aware chime file, enqueues playback.
  \item If \field{speak} is \texttt{"y"} (voice mode) and
    \field{vibe\_signals} is non-empty: selects a random phrase
    from the relevant quip pool, calls \texttt{vox unmute} to
    synthesize and play.  If \field{vibe\_signals} is empty
    (no bash commands have run yet), the Stop hook is a no-op
    even with voice enabled.
  \item For Stop hooks: returns
    \texttt{\{"decision": "block", "reason": phrase\}} to halt
    Claude's output while audio plays.
\end{enumerate}

\subsection{Session Lifecycle}

\begin{enumerate}
  \item \textbf{Session start.}  \hook{SessionStart} hook runs
    \texttt{session-start.sh}: deploys slash commands, cleans
    retired commands, auto-allows MCP tools.  No audio plays.
  \item \textbf{User enables voice.}  \cmd{unmute} or \cmd{vox}
    sets \field{notify} to \texttt{"y"} or \texttt{"c"}.  This
    is the first moment audio becomes possible.
  \item \textbf{Signal accumulation.}  Each \hook{PostToolUse}
    (Bash) event classifies stdout against regex patterns and
    appends a timestamped signal (e.g.,
    \texttt{tests-pass@10:45}) to \field{vibe\_signals} in
    config.
  \item \textbf{Mood inference.}  On Stop hooks, the signal
    history is deterministically mapped to ElevenLabs expressive
    tags (e.g., \texttt{[relieved] [satisfied]}) without any
    LLM call.
  \item \textbf{Session end.}  \hook{SessionEnd} hook speaks a
    farewell phrase and clears \field{vibe\_signals}.
\end{enumerate}


% ===================================================================
\section{Provider Architecture}
% ===================================================================

\subsection{Provider Protocol}

All providers implement \texttt{TTSProvider}, a Python
\texttt{Protocol} class extending \texttt{AudioProvider}.  Eleven
required members:

\begin{center}
\begin{tabular}{lp{8cm}}
\toprule
\textbf{Member} & \textbf{Signature and Purpose} \\
\midrule
\field{name}           & \texttt{str} --- Short identifier
  (e.g., \texttt{"elevenlabs"}) \\
\field{default\_voice} & \texttt{str} --- Fallback voice name \\
\field{supports\_expressive\_tags} & \texttt{bool} --- Whether the
  provider interprets \texttt{[bracketed]} tags in text \\
\texttt{synthesize}    & \texttt{(request, output\_path) -> AudioResult}
  --- Core synthesis \\
\texttt{resolve\_voice} & \texttt{(name, language?) -> str} --- Voice
  validation and normalization \\
\texttt{get\_default\_voice} & \texttt{(language) -> str} ---
  Language-aware default \\
\texttt{list\_voices}  & \texttt{(language?) -> list[str]} ---
  Available voices \\
\texttt{infer\_language\_from\_voice} & \texttt{(voice) -> str | None}
  --- Reverse lookup (best-effort) \\
\texttt{check\_health} & \texttt{() -> list[HealthCheck]} ---
  Provider-specific diagnostics \\
\texttt{generate\_audio} & \texttt{(text, voice?, ...) -> AudioResult}
  --- Convenience wrapper (inherited from \texttt{AudioProvider}) \\
\texttt{generate\_audios} & \texttt{(texts, voice?, ...) -> list[AudioResult]}
  --- Batch convenience wrapper \\
\bottomrule
\end{tabular}
\end{center}

No provider SDK imports appear outside the provider's own file.
Providers do import shared utilities from \texttt{core.py}
(\texttt{split\_text}, \texttt{stitch\_audio}) for chunking, but
the SDK libraries (\texttt{elevenlabs}, \texttt{openai},
\texttt{boto3}) are confined to their respective files.

\subsection{ElevenLabs}

\begin{center}
\begin{tabular}{ll}
\toprule
\textbf{Parameter} & \textbf{Value} \\
\midrule
Default model   & \texttt{eleven\_flash\_v2\_5} (low latency, 32 languages) \\
Expressive tags & Only on \texttt{eleven\_v3} \\
Output format   & MP3 44.1\,kHz 128\,kbps \\
Char limits     & 5\,000--40\,000 depending on model \\
Chunking        & Auto-split at limit, stitch with 0\,ms pause \\
Voice resolution & API fetch (cached globally), match by lowercase name
  or raw 20-char voice ID \\
Voice settings  & stability, similarity\_boost, style, speaker\_boost
  (all optional floats 0.0--1.0) \\
\bottomrule
\end{tabular}
\end{center}

Streaming: the SDK returns a generator of audio chunks.  Vox
iterates the generator to completion and writes the full file.
There is no incremental playback---the entire response is
buffered to disk before \texttt{enqueue()} is called.

\subsection{OpenAI}

\begin{center}
\begin{tabular}{ll}
\toprule
\textbf{Parameter} & \textbf{Value} \\
\midrule
Default model   & \texttt{tts-1} (\texttt{tts-1-hd} available) \\
Expressive tags & Not supported \\
Max chars/call  & 4\,096 (hard API limit) \\
Chunking        & Auto-split at 4\,096 chars, stitch \\
Voices          & 9 static: alloy, ash, coral, echo, fable, nova,
  onyx, sage, shimmer \\
Rate mapping    & CLI percentage $\to$ OpenAI speed float
  (0.25--4.0, direct linear: rate/100) \\
\bottomrule
\end{tabular}
\end{center}

\subsection{AWS Polly}

\begin{center}
\begin{tabular}{ll}
\toprule
\textbf{Parameter} & \textbf{Value} \\
\midrule
Output format   & MP3 \\
Engine preference & neural $>$ generative $>$ long-form $>$ standard \\
Voice resolution & API fetch (paginated \texttt{DescribeVoices}, cached) \\
Language mapping & ISO 639-1 $\to$ BCP 47 (e.g., \texttt{de} $\to$
  \texttt{de-DE}) \\
Auth            & AWS credential chain (\texttt{aws sts
  get-caller-identity}) \\
\bottomrule
\end{tabular}
\end{center}

Voice resolution bundles three Polly parameters
(\field{voice\_id}, \field{language\_code}, \field{engine}) into
a frozen \texttt{VoiceConfig} dataclass, so callers never need to
know Polly's multi-parameter API.

\subsection{Offline Fallbacks}

\begin{center}
\begin{tabular}{llp{6cm}}
\toprule
\textbf{Provider} & \textbf{Platform} & \textbf{Notes} \\
\midrule
\texttt{say}     & macOS & System \texttt{say} command.  Default voice:
  \texttt{samantha}.  Output: AIFF converted to MP3 via ffmpeg.
  No API key required.  No expressive tags. \\
\texttt{espeak}  & Linux & \texttt{espeak-ng} command.  Voice is the
  language code (e.g., \texttt{en}).  No API key required.  No
  expressive tags. \\
\bottomrule
\end{tabular}
\end{center}

These exist so \texttt{vox doctor} and basic functionality work
on machines with no cloud credentials.  Audio quality is
noticeably lower than cloud providers.

\subsection{Provider Selection}

Auto-detection follows a strict priority chain:

\begin{enumerate}
  \item \texttt{TTS\_PROVIDER} environment variable (explicit
    override).
  \item \texttt{ELEVENLABS\_API\_KEY} present $\to$
    \texttt{elevenlabs}.
  \item \texttt{OPENAI\_API\_KEY} present $\to$ \texttt{openai}.
  \item \texttt{aws sts get-caller-identity} succeeds $\to$
    \texttt{polly}.
  \item macOS with \texttt{say} binary $\to$ \texttt{say}.
  \item Linux with \texttt{espeak-ng} or \texttt{espeak} $\to$
    \texttt{espeak}.
  \item Fallback: \texttt{polly} with a warning (will fail at
    synthesis time if no credentials).
\end{enumerate}

Session-level overrides persist to the durable config
(\texttt{.punt-labs/vox/vox.md}) when set via MCP tools.  The
\texttt{get\_provider()} factory checks:
explicit parameter $\to$ session config provider $\to$
auto-detect.  Model resolution follows the same cascade:
explicit $\to$ session config (only if provider matches) $\to$
\texttt{TTS\_MODEL} env $\to$ provider default.

\textbf{What this does not do:} there is no failover.  If the
selected provider's API call fails, the error propagates.  There
is no automatic retry, no fallback to a cheaper provider, and no
circuit breaker.


% ===================================================================
\section{MCP Tool Surface}
\label{sec:mcp}
% ===================================================================

The MCP server registers as key \texttt{mic} and exposes eleven
tools: four speech/state tools (\tool{unmute}, \tool{record},
\tool{vibe}, \tool{who}), three notification tools (\tool{notify},
\tool{speak}, \tool{status}), and the four-tool music suite
(\tool{music}, \tool{music\_play}, \tool{music\_list},
\tool{music\_next}).  All tools return JSON strings.

\begin{center}
\begin{longtable}{lp{8.5cm}}
\toprule
\textbf{Tool} & \textbf{Signature and Purpose} \\
\midrule
\endfirsthead
\toprule
\textbf{Tool} & \textbf{Signature and Purpose} \\
\midrule
\endhead
\tool{unmute} & \texttt{(text?, voice?, language?, segments?,
  rate=90, pause\_ms=500, ephemeral=true, stability?, similarity?,
  style?, speaker\_boost?, vibe\_tags?, provider?, model?)} ---
  Synthesize and play audio.  \field{segments} is a list of
  \texttt{\{text, voice?, language?, vibe\_tags?\}} dicts for
  multi-voice sequential playback.  Returns immediately; synthesis
  and playback happen in a background daemon thread.  Persists
  provider, model, and vibe\_tags to session config. \\
\tool{record} & \texttt{(text?, voice?, language?, segments?,
  rate=90, pause\_ms=500, output\_path?, output\_dir?,
  stability?, similarity?, style?, speaker\_boost?)} --- Like
  \tool{unmute} but writes to persistent path (default
  \texttt{\~{}/Music/vox/}), no playback. \\
\tool{vibe}   & \texttt{(mood?, tags?, mode?)} --- Set session
  mood, expressive tags, or vibe mode
  (\texttt{auto}/\texttt{manual}/\texttt{off}). \\
\tool{who}    & \texttt{(language?)} --- List available voices:
  6 random featured voices with personality blurbs, plus full
  roster.  Filters by language if provided. \\
\tool{notify} & \texttt{(mode?, voice?)} --- Set notification
  mode: \texttt{y} (on-demand), \texttt{n} (off), \texttt{c}
  (continuous).  First enable auto-sets \field{speak} to
  \texttt{"y"}. \\
\tool{speak}  & \texttt{(mode?, voice?)} --- Set speech mode:
  \texttt{y} (voiced) or \texttt{n} (chimes only). \\
\tool{status} & \texttt{()} --- Returns full session state:
  provider, voice, notify, speak, vibe\_mode, vibe, vibe\_tags,
  vibe\_signals, music\_mode. \\
\tool{music} & \texttt{(mode, style?, name?, base\_prompt?, variations?)}
  --- Start (\texttt{"on"}) or stop (\texttt{"off"}) background music.
  When on, \texttt{voxd} builds (or reopens) a named, persisted
  \emph{Program}: it plays the first generated track at once and fills a
  pool of up to twelve distinct instrumental tracks in the background,
  auto-advancing as each ends and shuffle-rotating the full pool at zero
  credits.  The Program is daemon-owned and ownership-free --- any
  session drives it (DES-041), superseding the old session-ownership
  model.  A matching saved \field{name} reopens the pool with no
  generation. \\
\tool{music\_play} & \texttt{(name)} --- Reopen and play a saved Program
  by name.  No generation, no credits. \\
\tool{music\_list} & \texttt{()} --- List saved Programs with name, size,
  and date. \\
\tool{music\_next} & \texttt{()} --- Advance to the next track now: rotate
  the full pool (zero credits) or generate one if the pool is not yet
  full. \\
\bottomrule
\end{longtable}
\end{center}

\subsection{Background Synthesis Pattern}

Both \tool{unmute} and \tool{record} return immediately with
\emph{predicted} results.  The server computes deterministic output
paths (MD5 hash of text $\to$ 12-hex-char filename) and returns
metadata before synthesis begins.  A daemon thread performs the
actual provider API call and playback.

This means the agent sees the response in milliseconds, but audio
may not play for several seconds (depending on provider latency).
If synthesis fails in the background thread, the error is logged
but not reported to the agent.  The agent has no mechanism to
detect synthesis failure after the tool call returns.

\subsection{Transport}

The MCP server runs on stdio transport (\texttt{stdin}/\texttt{stdout}).
Logs go to stderr only.  The server is started by Claude Code via
\texttt{vox mcp} and communicates using the MCP JSON-RPC protocol.
The MCP server delegates synthesis and playback to
\texttt{voxd}, a persistent WebSocket daemon.  Each Claude Code
session spawns its own MCP server process, but all sessions
share a single \texttt{voxd} instance.


% ===================================================================
\section{Audio Pipeline}
% ===================================================================

\subsection{Text Splitting}

\texttt{split\_text(text, max\_chars)} guarantees every chunk is
$\leq$ \texttt{max\_chars}:

\begin{enumerate}
  \item If text fits: return as-is.
  \item Split at sentence boundaries (regex
    \texttt{(?<=[.!?])\textbackslash{}s+}).
  \item Accumulate sentences into chunks up to the limit.
  \item If a sentence exceeds the limit: split at word boundaries.
  \item If a word exceeds the limit: split at character boundaries.
\end{enumerate}

Whitespace between sentences is normalized to a single space.
This function is provider-agnostic and used by both ElevenLabs
(5\,000--40\,000 char models) and OpenAI (4\,096 char hard limit).

\subsection{Segment Processing}

Multi-segment requests (via the \field{segments} parameter) are
processed sequentially:

\begin{enumerate}
  \item Each segment resolves its own voice and language
    independently, with per-segment overrides falling back to
    top-level defaults.
  \item Per-segment \field{vibe\_tags} override the top-level
    default (only for providers that support expressive tags).
  \item Segments are synthesized to individual temp files.
  \item \texttt{stitch\_audio()} concatenates all segment files
    with \field{pause\_ms} silence between each (default 500\,ms).
  \item 150\,ms trailing silence is appended to prevent MP3 frame
    truncation artifacts.
\end{enumerate}

\subsection{Audio Stitching}

\texttt{stitch\_audio(segments, output\_path, pause\_ms)} uses
pydub's \texttt{AudioSegment}:

\begin{itemize}
  \item Loads each MP3 segment file.
  \item Inserts \texttt{AudioSegment.silent(duration=pause\_ms)}
    between segments.
  \item Appends 150\,ms trailing silence.
  \item Exports as MP3 to \field{output\_path}.
  \item Raises \texttt{FileNotFoundError} if any segment is
    missing; \texttt{ValueError} if the segment list is empty.
\end{itemize}

pydub delegates to ffmpeg for MP3 decode/encode.  This is the only
place in the system that requires ffmpeg at runtime (playback uses
\texttt{afplay} or \texttt{ffplay}, but stitching uses ffmpeg
via pydub).

\subsection{Playback Queue}

Playback is serialized across all concurrent callers via POSIX
file locking:

\begin{center}
\begin{tabular}{lp{8cm}}
\toprule
\textbf{Component} & \textbf{Details} \\
\midrule
Lock file      & \texttt{\~{}/.punt-labs/vox/playback.lock} \\
Lock type      & \texttt{fcntl.flock(fd, LOCK\_EX)} (blocking) \\
Timeout        & 120 seconds \\
Player (macOS) & \texttt{afplay} \\
Player (other) & \texttt{ffplay -nodisp -autoexit -loglevel quiet} \\
Pending dir    & \texttt{\~{}/.punt-labs/vox/pending/} \\
\bottomrule
\end{tabular}
\end{center}

{\emergencystretch=3em\relax
\texttt{enqueue(path)} copies the audio file to the pending
directory (preventing deletion races with ephemeral cleanup), then
spawns a fully detached subprocess that acquires the lock, plays,
and deletes the pending copy.  The parent returns immediately.\par}

\textbf{What this does not do:} there is no priority queue, no
cancellation, no skip-ahead.  If three audio files are enqueued,
they play in FIFO order.  There is no way to interrupt a playing
file or flush the queue.

\subsection{Background Music}
\label{sec:music}

Since v4.x, \texttt{voxd} generates, persists, and plays instrumental
background music as an ownership-free \emph{Program}: a named on-disk
pool of tracks with an explicit, formally-modelled lifecycle
(\texttt{docs/audio-programs.tex}; DES-041).  This generalised and
replaced an earlier session-owned, single-looping-track design --- the
superseded DES-031 (session ownership) and DES-033 (vibe-triggered
regeneration).  Music deliberately sidesteps the speech playback queue:

\begin{itemize}
  \item \textbf{Ownership-free, daemon-owned (DES-041).} \texttt{voxd}
    holds exactly one Program.  There is no \field{owner\_id} and no
    non-owner rejection --- any session, and the CLI, drive the same
    Program.  A single-writer \texttt{ControlChannel} serialises every
    mutation as a typed \texttt{ControlSignal}, so concurrent clients
    never race the pool, and \texttt{ProgramState} re-validates all
    sixteen state invariants by construction on every transition.
  \item \textbf{Pool, auto-advance, and rotation.} \texttt{music("on")}
    plays the first generated track immediately, then a background filler
    generates up to twelve distinct tracks (one \texttt{asyncio} task at a
    time).  Playback auto-advances as each track ends; once the pool is
    full, generation stops and playback shuffle-rotates the pool ---
    never repeating the just-played track --- at zero credits.  A
    resilient retry sub-machine (\texttt{retry\_machine.py}) absorbs
    transient generation failures without corrupting the pool
    (see the \texttt{RetryCapped}/\texttt{RetryExhausted} transitions in
    the Z model).
  \item \textbf{Persisted and replayable.} Each Program is a directory
    \texttt{\~{}/Music/vox/<name>/} of \texttt{NNN.mp3} parts plus a
    manifest; every track carries ID3 tags (artist \texttt{vox}, album =
    program name, title = the variation clause, genre = style, and track
    index), so pools drop straight into Music.app or Ubuntu players.
    Both the CLI and MCP can \texttt{list}, \texttt{play}, \texttt{loop},
    and address a specific part (\texttt{playlist:N}) of a saved
    Program --- consume-only replay with no generation and no credits.
  \item \textbf{Separate subprocess at reduced volume (DES-030).} Each
    track plays through its own player process (\texttt{afplay
    -{}-volume 0.3} on macOS, \texttt{ffplay -volume 30} on Linux),
    completely outside the speech mutex.  Speech and chimes overlay at
    full volume; the volume differential makes speech intelligible over
    the music with no runtime ducking or pausing.
  \item \textbf{LLM-authored prompts.} The agent authors a
    \field{base\_prompt} plus twelve genre-accurate \field{variations}
    (scaffolded by \texttt{music\_prompts.py}); \texttt{voxd} is a pipe to
    ElevenLabs and never decides what a genre sounds like.  With no
    prompts supplied it falls back to a bare stylistic default.
\end{itemize}

\subsection{Ephemeral Mode}

When \field{ephemeral=true} (the default for \tool{unmute}),
output is ephemeral: previous ephemeral files are cleaned up before
each synthesis.  In the current daemon-centric design \texttt{voxd}
owns ephemeral cleanup internally (writing under
\texttt{.punt-labs/vox/ephemeral/}); the \field{ephemeral} flag is
still accepted on the tool surface for callers.

Persistent output (via \tool{record} or \field{ephemeral=false})
goes to \texttt{\~{}/Music/vox/} (overridable via
\texttt{VOX\_OUTPUT\_DIR}) or a user-specified directory.  Generated
music Programs live under \texttt{\~{}/Music/vox/<name>/} --- one
directory per named Program (\S\ref{sec:music}).

\subsection{Format Handling}

All providers output MP3.  Cloud providers (ElevenLabs, OpenAI,
Polly) produce MP3 natively.  The \texttt{say} provider produces
AIFF via the macOS \texttt{say} command, then converts to MP3
via ffmpeg.  The \texttt{espeak} provider writes WAV via
\texttt{espeak-ng}, then converts to MP3 via ffmpeg.
This means all providers produce files that pydub can decode,
and multi-segment stitching works across all backends when
ffmpeg is available.


% ===================================================================
\section{CLI and Config}
% ===================================================================

\subsection{CLI Commands}

\begin{center}
\begin{longtable}{lp{8cm}}
\toprule
\textbf{Command} & \textbf{Purpose} \\
\midrule
\endfirsthead
\toprule
\textbf{Command} & \textbf{Purpose} \\
\midrule
\endhead
\texttt{vox unmute [TEXT]} & Synthesize and play.
  \texttt{-{}-from FILE} for batch input (JSON array).
  Ephemeral output to \texttt{.vox/}. \\
\texttt{vox record [TEXT]} & Synthesize and save.
  \texttt{-{}-output PATH} or \texttt{-{}-output-dir DIR}.
  No playback. \\
\texttt{vox vibe MOOD}     & Set session mood.  Special values:
  \texttt{auto} (infer from signals), \texttt{off} (disable). \\
\texttt{vox music [on|off]} & Start/stop background music.
  \texttt{-{}-style} (e.g.\ \texttt{techno}), \texttt{-{}-name} to save
  or replay a track. \\
\texttt{vox notify [y|n|c]} & Set notification mode.
  \texttt{y} = on-demand, \texttt{n} = off, \texttt{c} = continuous.
  Optional \texttt{-{}-voice}. \\
\texttt{vox speak [y|n]}   & \texttt{y} = voiced, \texttt{n} = chimes
  only. \\
\texttt{vox voice NAME}    & Set session voice. \\
\texttt{vox status}        & Show session state (provider, voice,
  mood, signals). \\
\texttt{vox doctor}        & Health checks: Python version,
  player binary, provider API keys, provider-specific checks. \\
\texttt{vox version}       & Print version string. \\
\texttt{vox mcp}           & Start MCP server on stdio. \\
\texttt{vox daemon <cmd>}  & Manage the \texttt{voxd} system service:
  \texttt{install}, \texttt{uninstall}, \texttt{status},
  \texttt{restart}. \\
\texttt{vox hook <event>}  & Dispatch hook handler (stop,
  post-bash, notification, pre-compact, user-prompt-submit,
  subagent-start, subagent-stop, session-end). \\
\bottomrule
\end{longtable}
\end{center}

Global flags: \texttt{-{}-json} (JSON output), \texttt{-{}-verbose}
(debug logging), \texttt{-{}-quiet} (suppress non-JSON output).
\texttt{-{}-verbose} and \texttt{-{}-quiet} are mutually exclusive.

\subsection{Config File Format}

Session state is a \emph{split} pair of YAML-frontmatter files under
\texttt{.punt-labs/vox/} in the repo root (DES-036).  Durable
preferences are tracked in git; ephemeral session state is gitignored.
Which file a field lives in is decided by the \texttt{DURABLE\_KEYS} /
\texttt{EPHEMERAL\_KEYS} frozensets in \modpath{config.py}.

\begin{lstlisting}
# .punt-labs/vox/vox.md  (durable, tracked)
---
notify: "y"
speak: "y"
voice: "matilda"
provider: "elevenlabs"
model: "eleven_v3"
vibe_mode: "auto"
---

# .punt-labs/vox/vox.local.md  (ephemeral, gitignored)
---
vibe: "focused coding session"
vibe_tags: "[calm] [focused]"
vibe_signals: "tests-pass@10:45,git-commit@10:47"
---
\end{lstlisting}

\begin{center}
\begin{tabular}{lllp{4.3cm}}
\toprule
\textbf{Field} & \textbf{File} & \textbf{Values} & \textbf{Purpose} \\
\midrule
\field{notify}       & \texttt{vox.md} & \texttt{y}, \texttt{n}, \texttt{c} &
  Notification mode \\
\field{speak}        & \texttt{vox.md} & \texttt{y}, \texttt{n} &
  Voice vs.\ chime mode \\
\field{voice}        & \texttt{vox.md} & string or null &
  Session voice name \\
\field{provider}     & \texttt{vox.md} & string or null &
  TTS provider override \\
\field{model}        & \texttt{vox.md} & string or null &
  TTS model override \\
\field{vibe\_mode}   & \texttt{vox.md} & \texttt{auto/manual/off} &
  Mood tracking mode \\
\field{vibe}         & \texttt{vox.local.md} & string or null &
  Human-readable mood \\
\field{vibe\_tags}   & \texttt{vox.local.md} & string or null &
  ElevenLabs expressive tags \\
\field{vibe\_signals} & \texttt{vox.local.md} & string or null &
  Timestamped signal history \\
\bottomrule
\end{tabular}
\end{center}

\texttt{ConfigStore} owns a config directory (default
\texttt{.punt-labs/vox/}); \texttt{find\_config\_dir()} walks up from
cwd looking for a directory containing \texttt{vox.md} or
\texttt{vox.local.md} (no legacy \texttt{.vox/} fallback).  Reads merge
the two files --- durable as the base layer, ephemeral overlaid but
accepting only \texttt{EPHEMERAL\_KEYS}.

\texttt{write\_field()} and \texttt{write\_fields()} validate keys
against \texttt{ALLOWED\_CONFIG\_KEYS} before writing and route each key
to the correct file; unknown keys raise \texttt{ValueError}, and values
containing newlines are rejected to protect the frontmatter.  If the
target file does not exist, a minimal file is created with YAML
frontmatter delimiters.

\subsection{Credential Management}

Provider API keys are read from environment variables:

\begin{center}
\begin{tabular}{ll}
\toprule
\textbf{Variable} & \textbf{Provider} \\
\midrule
\texttt{ELEVENLABS\_API\_KEY} & ElevenLabs (TTS + music) \\
\texttt{OPENAI\_API\_KEY}     & OpenAI \\
AWS credential chain           & Polly (via boto3) \\
\bottomrule
\end{tabular}
\end{center}

Because \texttt{voxd} runs as a detached service, it cannot inherit the
shell environment of the session that installed it.  At \texttt{vox
daemon install}, \texttt{KeysEnvWriter} snapshots the relevant provider
keys from the installing environment into
\texttt{\~{}/.punt-labs/vox/keys.env} (\texttt{chmod 0700} directory) so
the daemon can read them at startup.  There is no keychain integration
and no encrypted storage.  \texttt{vox doctor} checks for the presence
of these variables and reports their status.

\subsubsection*{Per-call API keys}

The CLI also accepts a \emph{per-call} key that overrides the ambient
credentials for a single invocation --- for billing a synthesis to a
specific account.  \texttt{ApiKeyResolver} resolves it from exactly one
of four mutually-exclusive sources, rejecting any combination:

\begin{enumerate}
  \item \texttt{-{}-api-key-file <path>} (mode checked; warns if broader
    than \texttt{0600}).
  \item \texttt{-{}-api-key-stdin} (one line piped in; refuses a tty).
  \item \texttt{VOX\_API\_KEY} environment variable.
  \item \texttt{-{}-api-key <value>} on argv (warns that argv is visible
    via \texttt{ps} and shell history).
\end{enumerate}

Absent all four, the call uses the ambient provider credentials.

\subsubsection*{Remote daemon auth}

When a client talks to a remote \texttt{voxd} (DES-037), the auth token
is the security boundary.  It resolves \texttt{explicit arg > VOXD\_TOKEN
env > serve.token file}, and access logs redact it.  \texttt{VOXD\_HOST}
and \texttt{VOXD\_PORT} select the remote endpoint (default
\texttt{127.0.0.1} and the local \texttt{serve.port}); \texttt{VOXD\_BIND}
sets the daemon's bind address server-side (default \texttt{127.0.0.1}).


% ===================================================================
\section{Performance}
% ===================================================================

\subsection{Latency Characteristics}

\begin{center}
\begin{tabular}{lrl}
\toprule
\textbf{Provider} & \textbf{Typical Latency} & \textbf{Notes} \\
\midrule
ElevenLabs (flash) & 200--800\,ms & Streaming API, buffered to disk \\
ElevenLabs (v3)    & 500--2\,000\,ms & Higher quality, higher latency \\
OpenAI (tts-1)     & 300--1\,000\,ms & No streaming, single HTTP call \\
OpenAI (tts-1-hd)  & 500--1\,500\,ms & Higher quality variant \\
Polly (neural)     & 200--600\,ms & AWS region-dependent \\
macOS \texttt{say} & 50--200\,ms & Local, no network \\
\texttt{espeak-ng} & 20--100\,ms & Local, no network \\
\bottomrule
\end{tabular}
\end{center}

These are synthesis latencies (time to first byte from provider).
Total time-to-audio includes stitching (10--50\,ms for pydub
operations) and playback queue wait (0\,ms if queue is empty,
unbounded if other audio is playing).

\subsection{Caching}

There is no intentional audio caching.  Deterministic filenames
(MD5 hash of text) mean that repeated calls for the same text
produce the same output path.  The MCP server's
\texttt{\_synthesize\_segments()} skips synthesis when the output
file already exists, which provides incidental caching for
non-ephemeral calls.  In ephemeral mode,
\texttt{clean\_ephemeral()} deletes prior MP3 files before each
call, so the skip rarely triggers.

Voice lists are cached globally per provider (populated on first
call, never refreshed within a process).  Config path resolution
is cached via \texttt{lru\_cache(maxsize=1)}.

\subsection{Streaming vs.\ Batch}

ElevenLabs returns a streaming response (chunk generator).  Vox
iterates the generator to completion and writes the full file
before playback.  This adds latency equal to the full synthesis
time but simplifies the pipeline: every provider produces a
complete file, and the playback system has a single code path.

\textbf{What this means:} for a 30-second audio clip, the user
waits for the full 30 seconds of synthesis before hearing anything.
True streaming playback (play while synthesizing) is not
implemented.


% ===================================================================
\section{Security}
% ===================================================================

\subsection{Threat Model}

Vox runs in the user's process space with the user's credentials.
The primary trust boundary is the MCP tool interface: Claude Code
can call any tool, but tools only perform synthesis and
configuration---no file deletion, no network access beyond
provider APIs, no shell execution.

\begin{center}
\begin{tabular}{lp{5.5cm}p{4cm}}
\toprule
\textbf{Threat} & \textbf{Description} & \textbf{Mitigation} \\
\midrule
API key exposure & Provider keys in env vars visible to any
  process in the session. & Standard env var practice.  No
  logging of key values.  \texttt{doctor} reports presence, not
  content. \\
Unintended audio & Audio plays without user consent, e.g.,
  during a meeting. & Default \field{notify=n}.  No audio until
  explicit opt-in.  \cmd{mute} immediately silences. \\
Cost accumulation & Rapid synthesis calls exhaust provider
  quotas. & No rate limiting in Vox.  Provider-side quotas are
  the only protection. \\
Pending file race & Concurrent processes write to
  \texttt{pending/} with colliding filenames. & Filenames include
  PID prefix.  \texttt{flock} serializes playback. \\
\bottomrule
\end{tabular}
\end{center}

\subsection{No Audio Until Opt-In}

This is the system's primary safety property.  On fresh install:

\begin{itemize}
  \item \texttt{.vox/config.md} does not exist.
  \item \field{notify} defaults to \texttt{"n"}.
  \item All hook handlers check \field{notify} first and return
    immediately if \texttt{"n"}.
  \item The \hook{SessionStart} hook deploys commands but
    produces no audio.
  \item No MCP tool produces audio unless the agent explicitly
    calls \tool{unmute}.
\end{itemize}

The user must take an affirmative action (\cmd{unmute},
\cmd{vox~y}, or the agent calling \tool{notify(mode="y")}) to
enable any audio output.

\subsection{Credential Isolation}

Each provider file is the only file that imports its SDK:

\begin{itemize}
  \item \modpath{providers/elevenlabs.py} --- only file importing
    \texttt{elevenlabs}
  \item \modpath{providers/openai.py} --- only file importing
    \texttt{openai}
  \item \modpath{providers/polly.py} --- only file importing
    \texttt{boto3}
\end{itemize}

API clients are instantiated inside provider constructors.  No
client objects are passed between modules.  No credentials are
logged at any level.


% ===================================================================
\section{Platform Integration}
% ===================================================================

\subsection{macOS}

\begin{itemize}
  \item \textbf{Playback.}  \texttt{afplay} is the primary
    player, available on all macOS installations.
    \texttt{ffplay} is the fallback (requires Homebrew ffmpeg).
  \item \textbf{Offline fallback.}  The \texttt{say} provider
    uses the system's built-in speech synthesis.  Voice names
    are macOS-specific (e.g., \texttt{samantha}, \texttt{alex}).
  \item \textbf{File locking.}  \texttt{fcntl.flock()} works
    correctly on APFS and HFS+.
  \item \textbf{Temp directory.}  \texttt{.envrc} sets
    \texttt{TMPDIR} to \texttt{.tmp/} in the project root,
    keeping temp files workspace-local.
\end{itemize}

\subsection{Linux}

\begin{itemize}
  \item \textbf{Playback.}  \texttt{ffplay} from the
    \texttt{ffmpeg} package.  No \texttt{afplay} equivalent.
    \texttt{resolve\_player()} raises \texttt{FileNotFoundError}
    if neither is available.
  \item \textbf{Offline fallback.}  The \texttt{espeak} provider
    wraps \texttt{espeak-ng} (preferred) or \texttt{espeak}.
  \item \textbf{File locking.}  \texttt{fcntl.flock()} works on
    ext4, btrfs, and most Linux filesystems.  NFS mounts may not
    honor \texttt{flock}; this is not handled.
  \item \textbf{No headless mode.}  Playback requires an audio
    device.  Running in a headless container or SSH session
    without PulseAudio/ALSA will cause \texttt{ffplay} to fail.
    There is no detection or graceful degradation for this case.
\end{itemize}


% ===================================================================
\section{Plugin and Hook Integration}
% ===================================================================

Vox integrates with Claude Code through a plugin manifest
(\texttt{plugin.json}) and a hook registration file
(\texttt{hooks.json}).  The plugin registers an MCP server
(key: \texttt{mic}) and nine lifecycle hooks.

\subsection{Hook Registration}

\begin{center}
\begin{longtable}{llp{5.5cm}}
\toprule
\textbf{Event} & \textbf{Script} & \textbf{Behavior} \\
\midrule
\endfirsthead
\toprule
\textbf{Event} & \textbf{Script} & \textbf{Behavior} \\
\midrule
\endhead
\hook{SessionStart}       & \texttt{session-start.sh} &
  Deploy slash commands, clean retired commands, auto-allow MCP
  tools.  No audio. \\
\hook{PostToolUse} (MCP)  & \texttt{suppress-output.sh} &
  Format MCP tool output as compact UI panel text. \\
\hook{PostToolUse} (Bash) & \texttt{signal.sh} &
  Classify stdout, append signal to
  \field{vibe\_signals}. \\
\hook{Stop}               & \texttt{notify.sh} &
  Block with speech phrase or play chime.  Returns
  \texttt{\{decision: "block"\}} to pause Claude while audio
  plays. \\
\hook{Notification}       & \texttt{notify-permission.sh} &
  Play chime or speak permission/idle phrase (async). \\
\hook{PreCompact}         & \texttt{pre-compact.sh} &
  Speak ``grabbing a snack'' phrase (continuous mode only,
  async). \\
\hook{UserPromptSubmit}   & \texttt{acknowledge.sh} &
  Speak acknowledgment phrase (continuous mode only, async). \\
\hook{SubagentStart/Stop} & \texttt{subagent.sh} &
  Announce subagent spawn/completion (continuous mode only,
  async). \\
\hook{SessionEnd}         & \texttt{farewell.sh} &
  Speak farewell, clear \field{vibe\_signals} (async). \\
\bottomrule
\end{longtable}
\end{center}

\subsection{Mute/Unmute Lifecycle}

Three notification states control audio behavior:

\begin{center}
\begin{tabular}{lp{4cm}p{4cm}}
\toprule
\textbf{State} & \textbf{Hooks Fire?} & \textbf{Audio?} \\
\midrule
\field{notify=n} & All hooks return immediately &
  No audio from hooks.  Agent can still call \tool{unmute}
  directly. \\
\field{notify=y} & Stop, Notification, SessionEnd &
  Chimes (\field{speak=n}) or speech (\field{speak=y}) on
  task completion and permission requests. \\
\field{notify=c} & All hooks fire &
  Continuous narration: acknowledgments, subagent
  announcements, pre-compact phrases, farewells. \\
\bottomrule
\end{tabular}
\end{center}

\field{speak} is orthogonal to \field{notify}:

\begin{itemize}
  \item \field{speak=y}: hooks synthesize speech via provider.
  \item \field{speak=n}: hooks play pre-recorded chime MP3 files
    from \texttt{assets/} (no provider API call).
\end{itemize}

\subsection{Signal Classification}

Signal classification uses a single canonical pattern table in
\texttt{hooks.py}.  The \hook{PostToolUse} (Bash) hook calls
\texttt{classify\_signal()} directly.  The session watcher in
\texttt{watcher.py} delegates to the same function via
\texttt{classify\_output()}, passing \texttt{exit\_code=None}.
The eight regex-based signal patterns are:

\begin{center}
\begin{tabular}{lp{7cm}}
\toprule
\textbf{Signal} & \textbf{Detection Patterns} \\
\midrule
\texttt{lint-fail}     & \texttt{Found [0-9]+ error} \\
\texttt{lint-pass}     & \texttt{All checks passed}, \texttt{0 errors} \\
\texttt{tests-pass}    & \texttt{[0-9]+ passed}, \texttt{tests? ok},
  \texttt{$\checkmark$.*passed} \\
\texttt{tests-fail}    & \texttt{FAILED}, \texttt{AssertionError},
  \texttt{ERRORS?\textbackslash{}b} \\
\texttt{merge-conflict} & \texttt{CONFLICT} \\
\texttt{git-push-ok}   & \texttt{Everything up-to-date},
  \texttt{->.*main} \\
\texttt{git-commit}    & \texttt{\^{}\textbackslash{}[.+\textbackslash{}] .+}, \texttt{\^{}create mode} \\
\texttt{pr-created}    & \texttt{pull/[0-9]+},
  \texttt{created pull request} \\
\bottomrule
\end{tabular}
\end{center}

If no pattern matches but exit code is non-zero, the signal is
\texttt{cmd-fail}.  Signals are appended as
\texttt{\{signal\}@\{HH:MM\}} to the comma-separated
\field{vibe\_signals} config field.

\subsection{Mood Inference}

On Stop hooks, accumulated signals are deterministically mapped
to ElevenLabs expressive tags.  No LLM call is involved:

\begin{center}
\begin{tabular}{lp{6cm}}
\toprule
\textbf{Signal Trajectory} & \textbf{Tags} \\
\midrule
Recovery (fail $\to$ pass) & \texttt{[relieved]} \\
Shipped, no failures       & \texttt{[satisfied]} \\
Shipped after failures     & \texttt{[relieved] [satisfied]} \\
Ended with failure, many fails & \texttt{[frustrated] [sighs]} \\
Many passes, no failures   & \texttt{[excited]} \\
Some passes                & \texttt{[calm]} \\
Default                    & \texttt{[calm]} \\
\bottomrule
\end{tabular}
\end{center}

This only affects audio when the provider supports expressive
tags (\texttt{eleven\_v3} model only).  All other providers
ignore the tags silently.


% ===================================================================
\section{Known Limitations}
% ===================================================================

\subsection{No Streaming Playback}

Audio is fully synthesized to disk before playback begins.  For
long texts, this means the user waits for the entire synthesis to
complete.  ElevenLabs returns a streaming response, but Vox
buffers it entirely.  True streaming playback (play the first
chunk while subsequent chunks synthesize) is not implemented.

\subsection{No Intentional Audio Caching}

The MCP server skips synthesis when the output file already exists
at the deterministic path, but ephemeral cleanup deletes prior
files before most calls.  There is no cache layer, no TTL, and
no cache invalidation.  For non-ephemeral use patterns with
repetitive text, the file-exists check provides incidental
reuse.  For ephemeral calls (the common path), every synthesis
hits the provider API.

\subsection{No Provider Failover}

If the selected provider fails (API error, rate limit,
credentials expired), the error propagates to the caller.  There
is no automatic fallback to a lower-tier provider, no retry with
exponential backoff, and no circuit breaker.

\subsection{No Playback Cancellation}

Once audio is enqueued, it plays to completion.  There is no skip,
pause, cancel, or flush-queue mechanism.  If three long audio
clips are enqueued, they play sequentially with no way to
interrupt.

\subsection{Background Failure Opacity}

The \tool{unmute} and \tool{record} MCP tools return predicted
results before synthesis completes.  If synthesis fails in the
background thread (provider error, network timeout, disk full),
the error is logged to stderr but not reported to the agent.  The
agent has no mechanism to detect or recover from synthesis
failures after the tool call returns.  The session watcher's
autonomous speech also fails silently (bare
\texttt{except Exception} with logging).

\subsection{Concurrent Config Writes}

\texttt{write\_fields()} in \texttt{config.py} performs a
read-then-write with no file-level lock.  The MCP server's
background synthesis thread and \hook{PostToolUse} hook handlers
can write \field{vibe\_signals} concurrently.  Both paths do
unguarded read-modify-write on the same file.  The race window
is narrow but real in continuous mode with fast bash commands.
The playback lock (\texttt{playback.lock}) governs audio
playback, not config writes.

\subsection{NFS File Locking}

The playback serialization relies on \texttt{fcntl.flock()}.  On
NFS mounts, \texttt{flock} may not provide mutual exclusion.  If
\texttt{\~{}/.punt-labs/vox/} is on an NFS filesystem, concurrent
playback from multiple machines could overlap.

\subsection{No Headless Audio Detection}

On Linux, if no audio output device is available (headless
container, SSH without PulseAudio), \texttt{ffplay} fails
silently or with an error.  Vox does not detect this condition
and does not warn the user.

\subsection{Voice List Staleness}

Provider voice lists are cached globally on first access and
never refreshed within a process lifetime.  If a voice is added
or removed on the provider side, the change is invisible until
the process restarts (next Claude Code session).


% ===================================================================
\section{Test Architecture}
% ===================================================================

\begin{center}
\begin{tabular}{lp{7cm}}
\toprule
\textbf{Test File} & \textbf{What It Tests} \\
\midrule
\texttt{test\_types.py}             & Domain type validation,
  \texttt{generate\_filename()} determinism \\
\texttt{test\_core.py}              & Text splitting edge cases,
  \texttt{TTSClient} batch/pair synthesis, stitching \\
\texttt{test\_config.py}            & Config read/write, worktree
  resolution, field validation \\
\texttt{test\_output.py}            & Output path resolution, env
  var handling \\
\texttt{test\_normalize.py}          & Text normalization for speech \\
\texttt{test\_cache.py}             & MP3 quip cache, content addressing \\
\texttt{test\_playback.py}          & \texttt{flock} serialization,
  \texttt{enqueue()} subprocess lifecycle \\
\texttt{test\_hooks.py}             & Signal classification, mood
  inference, hook handler dispatch \\
\texttt{test\_cli.py}               & CLI command smoke tests \\
\texttt{test\_server.py}            & MCP tool behavior with mocked
  providers \\
\texttt{test\_elevenlabs\_provider.py} & Voice resolution, chunking,
  voice settings, health checks \\
\texttt{test\_openai\_provider.py}  & Rate mapping, chunking at 4\,096
  chars \\
\texttt{test\_polly\_provider.py}   & VoiceConfig bundling, engine
  preference, Polly API mocks \\
\texttt{test\_say\_provider.py}     & macOS say command construction \\
\texttt{test\_espeak\_provider.py}  & espeak-ng command construction \\
\texttt{test\_client.py}            & \texttt{VoxClient}/\texttt{VoxClientSync}
  WebSocket client for \texttt{voxd} \\
\texttt{test\_service\_*.py}        & Service lifecycle, split per backend:
  \texttt{installer}, \texttt{launchd}, \texttt{systemd}, \texttt{process},
  \texttt{keys\_env} \\
\texttt{test\_signal.py}            & \texttt{Signal}/\texttt{SignalLog}
  classification and accumulation \\
\texttt{cli\_music.py} tests, \texttt{tests/programs/} & The client-side
  music command surface and the \texttt{voxd/programs/} Program suite:
  state-machine invariants (\texttt{test\_invariants},
  \texttt{test\_properties}), the single-writer control channel, the
  filler and retry sub-machines, rotation, the filesystem store and
  manifest, and the per-command handlers \\
\texttt{test\_api\_key\_resolver.py} & Four-source mutual-exclusion key
  resolution and permission warnings \\
\texttt{test\_paths.py}             & Per-user path resolution and
  \texttt{ensure\_user\_dirs} permissions \\
\texttt{test\_doctor.py}, \texttt{test\_daemon\_restarter.py} & Health
  checks and stale-daemon restart \\
\bottomrule
\end{tabular}
\end{center}

The table above is representative, not exhaustive: the suite mirrors
every source module, including the split domain types
(\texttt{test\_types\_synthesis.py}), the hook envelope/payload modules,
the output formatter, and each provider (including
\texttt{elevenlabs\_music} and \texttt{voice\_resolver}).

Quality gates: \texttt{ruff} (lint + format), \texttt{mypy} (strict),
\texttt{pyright} (strict), \texttt{pytest}, \texttt{shellcheck}
(all shell scripts in \texttt{hooks/} and \texttt{scripts/}).
Zero violations required.

Polly test fixtures use \texttt{AudioSegment.silent(duration=50)}
to generate valid MP3 bytes, because pydub delegates to ffmpeg
which rejects fake byte sequences.  Provider mocks use
\texttt{side\_effect=lambda} rather than \texttt{return\_value} to
produce fresh mock objects per call.

\end{document}
